\documentclass[11pt,a4paper]{article}

% ================== Packages ==================
\usepackage[margin=1in]{geometry}
\usepackage{amsmath, amssymb, amsfonts}
\usepackage{mathtools}
\usepackage{siunitx}
\usepackage{booktabs}
\usepackage{enumitem}
\usepackage{lmodern}
\usepackage[T1]{fontenc}
\usepackage{microtype}
\usepackage{graphicx}
\usepackage{listings}
\usepackage{xcolor}

% ============== Listings setup (code) =========
\lstdefinestyle{mystyle}{
  basicstyle=\ttfamily\small,
  keywordstyle=\color{blue!70!black}\bfseries,
  commentstyle=\color{green!40!black}\itshape,
  stringstyle=\color{red!50!black},
  showstringspaces=false,
  frame=single,
  framerule=0.3pt,
  rulecolor=\color{black!30},
  numbers=left,
  numberstyle=\tiny,
  numbersep=8pt,
  tabsize=2,
  breaklines=true
}
\lstset{style=mystyle, language=Python}

% ================ Title =======================
\title{Grid Search, Random Search, K\,-Fold Cross-Validation, dan Tracking Hasil \\}
\author{Taufiqurrahman}
\date{}

\begin{document}
\maketitle

\section*{Tujuan Pembelajaran}
\begin{itemize}[leftmargin=1.2em]
  \item Memahami \textbf{Grid Search} dan \textbf{Random Search} untuk tuning hyperparameter.
  \item Menerapkan \textbf{K\,-Fold Cross-Validation} (termasuk variasi stratified dan time-aware).
  \item Menyusun \textbf{skema tracking hasil} eksperimen yang rapi dan reprodusibel.
\end{itemize}

\section{Dasar: Hyperparameter Tuning}
Misalkan model bergantung pada vektor hyperparameter
\[
\boldsymbol{\lambda} = (\lambda_1,\lambda_2,\dots,\lambda_d)\in\Lambda,
\]
dan fungsi skor validasi rata-rata (mis.\ akurasi atau F1) adalah
\[
\hat{J}(\boldsymbol{\lambda}) = \frac{1}{K}\sum_{k=1}^{K} s^{(k)}(\boldsymbol{\lambda}),
\]
dengan $s^{(k)}$ adalah skor pada lipatan (\emph{fold}) ke-$k$.

Tujuan tuning:
\[
\boldsymbol{\lambda}^\star \in \arg\max_{\boldsymbol{\lambda}\in\Lambda}\ \hat{J}(\boldsymbol{\lambda}).
\]

\subsection{Grid Search}
\paragraph{Ide.} Menyapu kisi (\emph{grid}) semua kombinasi kandidat setiap hyperparameter.

\paragraph{Keunggulan.}
\begin{itemize}[leftmargin=1.2em]
  \item Implementasi sangat sederhana.
  \item Reproducible: ruang yang dievaluasi deterministik.
\end{itemize}

\paragraph{Keterbatasan.}
\begin{itemize}[leftmargin=1.2em]
  \item \textbf{Eksplosi kombinasi} saat dimensi meningkat; biaya komputasi cepat membesar.
  \item Distribusi evaluasi tidak adaptif terhadap sensitivitas tiap dimensi.
\end{itemize}

\paragraph{Kompleksitas.}
Jika tiap dimensi memiliki $m_i$ kandidat, total evaluasi
\[
M_{\text{grid}} = \prod_{i=1}^{d} m_i.
\]

\subsection{Random Search}
\paragraph{Ide.} Mengambil sampel acak $\boldsymbol{\lambda}^{(1)},\dots,\boldsymbol{\lambda}^{(N)}$ dari distribusi prior pada ruang hyperparameter.

\paragraph{Keunggulan.}
\begin{itemize}[leftmargin=1.2em]
  \item \textbf{Efisien di dimensi tinggi}: menutup kemungkinan pemborosan di dimensi yang tidak penting.
  \item Probabilistik: peluang menjangkau daerah baik meningkat dengan jumlah sampel.
\end{itemize}

\paragraph{Intuisi probabilistik.} Misalkan proporsi volume ruang yang ``bagus'' adalah $p\in(0,1]$. Probabilitas \emph{gagal} menemukan daerah bagus setelah $N$ sampel acak adalah
\[
P(\text{gagal}) = (1-p)^N,
\]
sehingga \emph{berhasil} minimal sekali adalah
\[
P(\text{berhasil}) = 1 - (1-p)^N.
\]
Saat $N$ naik, peluang menemukan setelan baik meningkat cepat bahkan untuk $p$ kecil.

\paragraph{Keterbatasan.}
\begin{itemize}[leftmargin=1.2em]
  \item Hasil bervariasi antar-ulang; butuh \emph{seed} untuk reproduksibilitas.
  \item Tanpa pengetahuan awal, distribusi sampel bisa kurang terarah (namun dapat diperbaiki dengan \emph{log-uniform} untuk skala lebar).
\end{itemize}

\subsection{Pemilihan Ruang Pencarian}
\begin{itemize}[leftmargin=1.2em]
  \item \textbf{Skala log} untuk parameter positif yang rentangnya lebar, mis.\ $C\in[10^{-3},10^{3}]$:
  \[
  \log_{10} C \sim \mathcal{U}(-3, 3).
  \]
  \item \textbf{Distribusi diskret} untuk integer seperti jumlah tetangga $k\in\{1,3,5,7,\dots\}$.
  \item \textbf{Kategori} untuk pilihan kernel, fungsi aktivasi, atau strategi penyeimbang kelas.
\end{itemize}

\section{K\,-Fold Cross-Validation (CV)}
\subsection{Definisi}
Bagi data menjadi $K$ lipatan berukuran hampir sama: $\mathcal{D} = \cup_{k=1}^K \mathcal{D}_k$ dan saling lepas. Untuk setiap $k$:
\[
\mathcal{D}_{\text{train}}^{(k)} = \mathcal{D}\setminus \mathcal{D}_k,\quad
\mathcal{D}_{\text{val}}^{(k)} = \mathcal{D}_k.
\]
Latih model pada $\mathcal{D}_{\text{train}}^{(k)}$, evaluasi pada $\mathcal{D}_{\text{val}}^{(k)}$ menghasilkan skor $s^{(k)}$. Estimator performa:
\[
\hat{\mu} = \frac{1}{K}\sum_{k=1}^{K} s^{(k)},\qquad
\hat{\sigma} = \sqrt{\frac{1}{K-1}\sum_{k=1}^{K}\big(s^{(k)}-\hat{\mu}\big)^2}.
\]

\subsection{Stratified K\,-Fold}
Untuk klasifikasi, menjaga proporsi kelas serupa di tiap lipatan mengurangi variansi estimator:
\[
\text{Proporsi kelas di } \mathcal{D}_k \approx \text{proporsi global}.
\]

\subsection{Time-Aware (Time Series CV)}
Jika data berurutan waktu, \textbf{jangan acak}:
\[
\begin{aligned}
\mathcal{D}_{\text{train}}^{(1)} &= [1 : T_1], \quad \mathcal{D}_{\text{val}}^{(1)} = (T_1 : T_2],\\
\mathcal{D}_{\text{train}}^{(2)} &= [1 : T_2], \quad \mathcal{D}_{\text{val}}^{(2)} = (T_2 : T_3], \ \text{dst.}
\end{aligned}
\]
Meniru situasi prediksi masa depan, mencegah \emph{leakage} informasi dari masa depan ke masa lalu.

\subsection{Ukuran $K$ dan Dampaknya}
\begin{itemize}[leftmargin=1.2em]
  \item $K$ kecil (mis.\ 3--5): latih lebih cepat, variansi estimator lebih tinggi.
  \item $K$ lebih besar (mis.\ 10): latih lebih lama, variansi lebih rendah.
  \item \emph{Leave-One-Out} (LOOCV, $K=n$): bias rendah, variansi tinggi, mahal komputasi.
\end{itemize}

\subsection{Nested CV (Singkat)}
Untuk menghindari \emph{optimism bias} pada tuning, gunakan CV \emph{luar} untuk mengukur performa akhir dan CV \emph{dalam} untuk memilih hyperparameter. Estimator akhir:
\[
\hat{J}_{\text{nested}} = \frac{1}{K_{\text{outer}}}\sum_{o=1}^{K_{\text{outer}}}\ \max_{\boldsymbol{\lambda}\in\Lambda}\ \frac{1}{K_{\text{inner}}}\sum_{i=1}^{K_{\text{inner}}} s^{(o,i)}(\boldsymbol{\lambda}).
\]

\section{Praktik Anti-Leakage dalam CV}
\begin{itemize}[leftmargin=1.2em]
  \item \textbf{Semua} transformasi berbasis data (\emph{imputation}, \emph{scaling}, \emph{encoding}, \emph{feature selection}) harus dipelajari \emph{hanya} dari \(\mathcal{D}_{\text{train}}^{(k)}\) dan kemudian diaplikasikan ke \(\mathcal{D}_{\text{val}}^{(k)}\).
  \item Gunakan \textbf{pipeline} agar urutan transformasi konsisten dan aman.
  \item Pada time series, pastikan \textbf{cutoff waktu} benar dan tidak mencampur masa depan ke masa lalu.
\end{itemize}

\section{Grid vs Random: Rangkuman Praktis}
\begin{center}
\begin{tabular}{@{}lcc@{}}
\toprule
Aspek & Grid Search & Random Search \\
\midrule
Sederhana & Sangat mudah & Mudah \\
Skalabilitas dimensi & Buruk (eksplosi kombinasi) & Baik \\
Eksplorasi awal & Kurang fleksibel & Sangat cocok \\
Reproducible & Ya (deterministik) & Ya (dengan seed) \\
Menemukan interaksi halus & Baik pada ruang kecil & Baik dengan $N$ memadai \\
\bottomrule
\end{tabular}
\end{center}

\section{Tracking Hasil Eksperimen}
\subsection{Mengapa Perlu Tracking?}
\begin{itemize}[leftmargin=1.2em]
  \item Menjaga \textbf{reprodusibilitas}: tahu persis konfigurasi yang menghasilkan skor.
  \item Memungkinkan \textbf{perbandingan adil}: metrik, data split, versi paket.
  \item Memudahkan \textbf{audit dan kolaborasi}.
\end{itemize}

\subsection{Apa yang Harus Dicatat?}
\begin{itemize}[leftmargin=1.2em]
  \item \textbf{Metadata}: waktu, \emph{experiment\_id}, \emph{run\_id}, \emph{seed}, hostname.
  \item \textbf{Data}: sumber dataset, ukuran train/val/test, skema fitur, \emph{hash} data (opsional).
  \item \textbf{Konfigurasi}: nama model, ruang hyperparameter, nilai terpilih.
  \item \textbf{Prosedur}: jenis CV ($K$, stratified/time-aware), \emph{pipeline} transformasi.
  \item \textbf{Metrik}: mean $\hat{\mu}$, std $\hat{\sigma}$, confusion matrix, ROC-AUC, dsb.
  \item \textbf{Artefak}: model terserialisasi, \emph{scaler/encoder}, grafik penting.
  \item \textbf{Versi}: commit git, versi paket Python, OS.
\end{itemize}

\subsection{Format Pencatatan Sederhana}
\paragraph{CSV / JSON Log.} Satu baris per \emph{run}:
\[
\texttt{run\_id, time, model, space, method, K, seed, mean\_score, std\_score, best\_params, notes}
\]
Contoh (JSON singkat):
\begin{lstlisting}
{
  "run_id": "2025-12-01_14-10-22",
  "model": "LogisticRegression",
  "search": "random",
  "space": {"C": "loguniform[1e-3, 1e3]", "penalty": ["l2"]},
  "cv": {"type": "stratified", "k": 5, "seed": 42},
  "score": {"metric": "f1_macro", "mean": 0.842, "std": 0.011},
  "best_params": {"C": 2.31, "penalty": "l2"},
  "data": {"train": 12000, "val_total": 3000, "features": 512},
  "env": {"python": "3.12", "sklearn": "1.5.2"},
  "git": {"commit": "a1b2c3d"},
  "notes": "ngram=(1,2), min_df=3"
}
\end{lstlisting}

\paragraph{Struktur Direktori Artefak.}
\begin{lstlisting}
experiments/
├─ 2025-12-01_14-10-22/           # run_id
│  ├─ config.json
│  ├─ metrics.json
│  ├─ model.pkl                   # model terlatih
│  ├─ pipeline.pkl                # full pipeline (preprocess + model)
│  ├─ cm.png                      # confusion matrix
│  └─ requirements-freeze.txt     # daftar versi paket
\end{lstlisting}

\subsection{Checklist Reprodusibilitas}
\begin{itemize}[leftmargin=1.2em]
  \item Tetapkan \textbf{random seed} untuk semua komponen yang relevan.
  \item Simpan \textbf{split} (indeks fold) atau simpan \textbf{generator} CV.
  \item \textbf{Freeze} versi paket pada saat \emph{best run}.
  \item Simpan \textbf{pipeline lengkap} (anti-leakage) dan \textbf{konfigurasi}.
\end{itemize}

\section{Contoh Pseudocode (sklearn)}
\subsection{Grid Search + Stratified K\,-Fold}
\begin{lstlisting}
from sklearn.model_selection import StratifiedKFold, GridSearchCV
from sklearn.pipeline import Pipeline
from sklearn.preprocessing import StandardScaler
from sklearn.linear_model import LogisticRegression

pipe = Pipeline([
  ("sc", StandardScaler(with_mean=False)),  # contoh bila fitur sparse: with_mean=False
  ("clf", LogisticRegression(max_iter=1000))
])

param_grid = {
  "clf__C": [0.01, 0.1, 1, 10, 100],
  "clf__penalty": ["l2"],
}

cv = StratifiedKFold(n_splits=5, shuffle=True, random_state=42)
search = GridSearchCV(pipe, param_grid, scoring="f1_macro", cv=cv, n_jobs=-1)
search.fit(X, y)

best = search.best_params_
mean = search.best_score_
\end{lstlisting}

\subsection{Random Search (log-uniform) + Tracking Ringkas}
\begin{lstlisting}
import numpy as np, json, time
from scipy.stats import loguniform
from sklearn.model_selection import StratifiedKFold, RandomizedSearchCV

param_dist = {
  "clf__C": loguniform(1e-3, 1e3),
}
cv = StratifiedKFold(n_splits=5, shuffle=True, random_state=42)
rnd = RandomizedSearchCV(pipe, param_distributions=param_dist,
                         n_iter=40, scoring="f1_macro", cv=cv, n_jobs=-1,
                         random_state=42)
rnd.fit(X, y)

run_id = time.strftime("%Y-%m-%d_%H-%M-%S")
log = {
  "run_id": run_id,
  "search": "random",
  "best_params": rnd.best_params_,
  "best_score": float(rnd.best_score_)
}
with open(f"experiments/{run_id}_metrics.json", "w") as f:
  json.dump(log, f, indent=2)
\end{lstlisting}

\section{Strategi Praktis}
\begin{itemize}[leftmargin=1.2em]
  \item Mulai dengan \textbf{random search} untuk eksplorasi cepat; lanjutkan \textbf{grid lokal} di sekitar rentang menjanjikan.
  \item Gunakan \textbf{skala log} untuk hyperparameter sensitif skala (mis.\ $C$, learning rate).
  \item Pastikan \textbf{CV sesuai data}: stratified untuk klasifikasi; time-aware untuk data berurut.
  \item Terapkan \textbf{pipeline} agar aman dari \emph{leakage} dan mudah diserialisasi.
  \item \textbf{Track} semua run penting; simpan artefak agar mudah \emph{roll-back}.
\end{itemize}

\section{Ringkasan}
\begin{itemize}[leftmargin=1.2em]
  \item \textbf{Grid} cocok untuk ruang kecil dan kebutuhan deterministik.
  \item \textbf{Random} efektif di ruang besar/tinggi dimensi; peluang temuan baik meningkat dengan jumlah sampel.
  \item \textbf{K\,-Fold} memberikan estimasi performa yang lebih stabil dibanding satu split; pilih variasi yang sesuai (stratified / time-aware).
  \item \textbf{Tracking hasil} adalah fondasi reprodusibilitas dan kolaborasi yang sehat.
\end{itemize}

\end{document}
